\documentclass[11pt, oneside]{book}

\usepackage[spanish, catalan, english]{babel}
\usepackage[utf8]{inputenc}
\usepackage[T1]{fontenc} % from http://www.monperrus.net/martin/producing+searchable+and+copyable+pdf+files+with+accents+using+latex-pdflatex
\usepackage{graphicx}
\usepackage{caption}
\usepackage{subcaption}
\usepackage{multirow}
\usepackage{array}
\usepackage{color, colortbl}
\usepackage{hhline}
\usepackage{eurosym}
\usepackage{parskip}
\usepackage[labelfont=bf]{caption} %labelfont is to make the Figure X.X: in bold
\usepackage[titletoc,toc,page]{appendix}
\usepackage{nameref}
\usepackage[hidelinks]{hyperref}
\usepackage{bookmark}
\usepackage{listings}
\usepackage[algo2e, ruled]{algorithm2e}
\usepackage{pdflscape}
\usepackage{titlesec}
\usepackage{lipsum} %FIXME
\usepackage[xindy, toc]{glossaries}
\usepackage{multicol}
% Tabular setting
\usepackage[table,xcdraw]{xcolor}
\newcolumntype{M}[1]{>{\centering\arraybackslash}m{#1}}
\newcolumntype{N}{@{}m{0pt}@{}}

% Todo notes
\newcommand\note[1]{\textcolor{red}{#1}}

% Footnote
\usepackage{footnote}
\makesavenoteenv{tabular}
\makesavenoteenv{table}

%So that subsubsections appear in the TOC and are numbered
\setcounter{secnumdepth}{3}
\setcounter{tocdepth}{3}

%http://tex.stackexchange.com/questions/11233/distance-between-chapter-title-and-text
\titleformat{\chapter}[hang]{\Huge\bfseries}{\thechapter}{15pt}{\Huge\bfseries}
\titlespacing{\chapter}{0pt}{-50pt}{20pt}

\selectlanguage{english}
\def\Tiny{\fontsize{5pt}{5pt} \selectfont}

\hypersetup
{
    %pdfborder={0 0 0 [0 0]}, % Treu tots bordes dels links
    %kidelinks,               % Treu els bordes dels links
    linkbordercolor={1 1 1},  % Bordes dels links de dins del doc en blanc
    %urlbordercolor={1 1 1},   % Bordes dels links de les url en blanc
    %pdfborder={1 1 0.45},     % Bordes però més prims (0.45)
        % colorlinks=true,          % Coloreja els links (anula linkbordercolor)
        % linkcolor=black,          % Els links en negre com la resta de text
    % Pdf options:
    pdfstartview=FitH,
    pdfauthor={Pablo Pardo Garcia},
    pdfsubject={Master Thesis},
    pdftitle={Automatic Age Estimation in Still Images},
    pdfkeywords={gesture recognition, NAO, wifibot, robotics, Msc Thesis},
}

%Commands for references
\newcommand{\myref}[2]{\hyperref[#2]{#1~\ref*{#2}}}
\newcommand{\figref}[1]{\hyperref[#1]{Figure~\ref*{#1}}}
\newcommand{\algref}[1]{\hyperref[#1]{Algorithm~\ref*{#1}}}
\newcommand{\coderef}[1]{\hyperref[#1]{Code~\ref*{#1}}}
\newcommand{\chapterref}[1]{\hyperref[#1]{\ref*{#1}~\nameref*{#1}}}
\newcommand{\annexref}[1]{\hyperref[#1]{Appendix \ref*{#1}~\nameref*{#1}}}

%%%%%%%%%%%%%%%%%%%%%%%%%%%%%%%%%%%%%%%%%%%%%%%%%%%%%%%%%%%%%%%%%%%%%%%
\DontPrintSemicolon
% To make listings and algorithms in the same table of contents.
% from http://tex.stackexchange.com/questions/115023/combine-listoflistings-and-listofalgorithms-into-one-list
\makeatletter
\AtBeginDocument{
    \let\c@algocf\c@lstlisting % Say algorithm2e to use the counter of listings
    \renewcommand{\algocf@list}{lol} % Say algorithm2e to write the material to the listings file lol
    \renewcommand*\l@algocf{\@dottedtocline{1}{1.5em}{2.3em}} % Make the indent of the entries of algorithm and listings equal
}
\makeatother
%%%%%%%%%%%%%%%%%%%%%%%%%%%%%%%%%%%%%%%%%%%%%%%%%%%%%%%%%%%%%%%%%%%%%%%

\lstset{
    breaklines=true,
    captionpos=b,
    basicstyle=\footnotesize \ttfamily,
    frame=single
}
\renewcommand{\lstlistlistingname}{List of Algorithms}


%\title{\textbf{Automatic Age Estimation in Still Images}}
%\author{Pablo Pardo Garcia \\ \\ \\
%       {\small \emph{Supervisor:} Sergio Escalera Guerrero} \\ 
%       {\small \emph{Cosupervisor:} Marc Oliu} }
%\date{February 2015}

%\newcommand{\HRule}{\rule{\linewidth}{0.5mm}}

%% Glossary
\makeglossaries
%http://en.wikibooks.org/wiki/LaTeX/Glossary
%https://www.sharelatex.com/learn/Glossaries#Terms


\newacronym{cnn}{CNN}{Convolutional Neural Network}
\newacronym{ann}{ANN}{Artificial Neural Network}
\newacronym{mlp}{MLP}{Multi Layer Perceptron}
\newacronym{knn}{kNN}{k-Nearest Neighbours}
\newacronym{svm}{SVM}{Support Vector Machines}
\newacronym{svr}{SVR}{Support Vector Regressor}
\newacronym{bif}{BIF}{Biologically Inspired Features}
\newacronym{aam}{AAM}{Active Appearance Models}
\newacronym{ages}{AGES}{AGing pattErn Subspace}
\newacronym{pca}{PCA}{Principal Component Analysis}
\newacronym{lbp}{LBP}{Local Binary Patterns}
\newacronym{mae}{MAE}{Mean Absolute Error}
\newacronym{cs}{CS}{Cumulative Score}
\newacronym{rbf}{RBF}{Radial Basis Function}

%\newglossaryentry{glossaryexample}
%{
%    name={},
%    description={},
%    plural={}
%}

\begin{document}

%Title page
%\maketitle
%\thispagestyle{empty}
\setcounter{page}{3}
\thispagestyle{empty}
\vspace*{-2cm}

\hbox{
\includegraphics[width=0.7cm]{figures/logoupc.eps}
\includegraphics[width=0.5cm]{figures/logoub.eps}
\includegraphics[width=1cm]{figures/logourv.eps}
\Large \bf Master in Artificial Intelligence (UPC-UB-URV)}
\hrule

%\vfill
%\vfill
\bigskip\bigskip\bigskip

\begin{center}

{\LARGE Master of Science Thesis}

\bigskip\bigskip\bigskip\bigskip\bigskip

\textbf{\huge \bf Automatic Age Estimation in Still Images}

\bigskip\bigskip\bigskip\bigskip\bigskip

{\LARGE \bf Pablo Pardo Garcia}

\vspace*{5cm}

\begin{multicols}{2}
{\large Supervisor:}

\medskip\medskip\smallskip

{\Large\bf Sergio Escalera Guerrero}

\medskip\medskip

{\Large Dept. Matemàtica Aplicada i Anàlisi}

\medskip

{\Large University of Barcelona}

%{\Large Master in Artificial Intelligence}

\bigskip\bigskip\bigskip

{\large Co-supervisor:}

\medskip\medskip\smallskip

{\Large\bf Marc Oliu Simón}

\medskip\medskip

{\Large Dept. Matemàtica Aplicada i Anàlisi}

\medskip

{\Large University of Barcelona}
\end{multicols}

\medskip\medskip\medskip\medskip\medskip

% This has obviously to be changed
{\Large Barcelona, 24th April 2015}

\end{center}

%\clearpage

%\thispagestyle{empty}
%\ 
%\clearpage

\chapter*{} % empty chapter to have a blank page after the title
\thispagestyle{empty}
\chapter*{Acknowledgements}
\thispagestyle{empty}
To write some time in the future...

\pagenumbering{arabic}
\setcounter{page}{0}
\clearpage

%Table of contents
\newpage
\pagenumbering{roman} % Roman numerals
\chapter*{Abstract}

Age Estimation tackles the problem of automatically predicting the age of people from visual data (still images, video data, depth maps, etc.). 
This is a very challenging problem because ageing is affected by many factors \cite{shephard1997aging} such as bad habits, accidents, make up or facial hair.

This works focus on three different age definitions, \textit{real age} as the actual age (time elapsed since the person was born), \textit{apparent age} as the perceived age from humans from the visual appearance and \textit{estimated age} as the predicted age by a machine from the visual appearance.

In this work, a new age estimation face image database is presented containing for first time in the literature real and apparent age annotations. A study comparing the estimation of both type of ages was done by proposing two age estimation methods using the state of the art techniques, one based on \acrfull{bif} and the other based on deep learning.


As a result of this work, an international challenge is being organized and the results will be presented in the ICCV conference 2015 (under revision). \selectlanguage{english}
\tableofcontents
\listoftables
\listoffigures
\lstlistoflistings
\newpage
\pagenumbering{arabic} % Normal numerals

%Document
\chapter{Introduction} \label{chap:introduction}

Computer Vision is a very important field in Artificial Intelligence that exists since the 1960s when digital image processing by computers became possible. Since the beginning of artificial vision, facial analysis has been a major interest in the research community (not just in Computer Vision but in other scientific areas such as biology \cite{bhl24064}, psychology \cite{ekm02}, neuroscience \cite{freiwald2009face} and sociology \cite{kemper1978social}) because its difficulty and its applications. Some of its applications include automatic detection of facial expressions \cite{cohen2003facial}, face detection \cite{hsu2002face}, face recognition \cite{wright2009robust}, face verification \cite{taigman2014deepface} and automatic estimation of age \cite{4359348}, gender \cite{alexandre2010gender} and ethnicity \cite{hosoi2004ethnicity}.

Age estimation is a field within the facial analysis area in Computer Vision that tackles the problem of automatically predicting the age of people from visual data (still images, video data, depth maps, etc.). One of the main issues that the age estimation problem has is that there are many factors that influence human perception of age, some factors affect the aging of a person \cite{shephard1997aging}, such as smoking, drinking alcohol, doing sports, alimentation, etc. and others affect the face appealing such as scars, plastic surgery, make-up, facial hair, etc.

Some definitions should be established beforehand regarding the concept of human age:
\begin{itemize}
	\item \textit{Real age}: The actual age (number of years passed since the person was born).
	\item \textit{Apparent age}: Perceived age from humans from the visual appearance. 
	\item \textit{Estimated age}: The predicted age by a machine from the visual appearance.
\end{itemize}

\section{Goals}
This work aims to study the differences (if any) of automatic age estimation from real age labels and apparent age labels. Given that does not exist any face image dataset with this two label annotations, a database with such requisites has been created. In order to do so, a web-based application was developed using the \textit{Facebook API} to facilitate a collaborative and competitive collection of face images.

As a consequence of this work, the first database in the literature containing real age and apparent age annotations for the face images was created and analysed. This database will allow researchers to tackle a different and new sub-problem of age estimation, \textit{Apparent Age Estimation}. The most important methods in the state of the art are evaluated over our proposed database in this work.

In \textit{Real Age Estimation} other external factors such as time evolution, habits and surgeries have to be taken into account. However \textit{Apparent Age Estimation} is based purely in the perception field.


\section{Age Challenge}

Given the innovative aspect of the database, the HuPBA (Human Pose and Behaviour Analysis) research group \cite{HuPBA} together with the company ChaLearn \cite{ChaLearn} are going to prepare an international challenge competition later this year and present the results in a workshop organized by the team in the ICCV conference 2015 edition (under revision) within the ChaLearn Looking at People series \cite{LaP} \cite{SergioEscalera2014} \cite{Escalera:2013:MGR:2522848.2532595} \cite{conf/icmi/EscaleraGBRGAESASBS13}. The challenge and the workshop will be sponsored by companies such as Microsoft Research\footnote{\url{http://research.microsoft.com/en-us/}}, Amazon\footnote{\url{http://www.amazon.com/}}, IAPR\footnote{\url{http://www.iapr.org/}} among others.

The challenge pretends to establish a State of the Art techniques for Apparent Age Estimation and compare the methods and the results with the Real Age Estimation State of the Art.

Initial results regarding the database and the methods presented in this work has being published in the IJCNN.



\chapter{Background} \label{chap:background}

\section{Biological Inspired Features}

\section{Convolutional Neural Networks}

\chapter{State of the art} \label{chap:sota}

Age estimation has historically been one of the most challenging problems within the field of facial analysis \cite{5406526}\cite{han:age}. Despite its multiple applications in many different areas there are relatively few publications compared to other topics in facial analysis. This difficulty is due to many factors: 
\begin{itemize}
	\item Depending on the application scenario, the age estimation problem can be taken as a multiclass classification problem or a regression problem.
	\item Large databases are difficult to collect, especially series of chronological image from the same individuals.
	\item The factors that affect the ageing process are uncontrollable and person specific \cite{4284917}\cite{4359348}\cite{1709980}.
\end{itemize}

The age estimation problem has generally two stages or blocks. The first one is the age representation in the images and the second one the learning algorithm. Several techniques have been published in order to deal with these two stages \cite{5406526}. In the Table \ref{tab:age-methods} there is a summary with some of the most important age estimation methods.

\section{Age Representation}\label{sec:ageRepre}

Age Representation, i.e. the extraction of the features that will represent the ageing patterns, is a very important step in the age estimation process. A good age representation will contain enough variation of the data to express the full complexity of the problem. There are many ways in the literature to represent ageing factors from an image. The most influential ones are described next.

\subsection{Anthropomorphic Models}\label{subsec:antropo}
The first known work on age estimation from facial images was done by Y. Kwon and N. Lobo \cite{Kwon:1999:ACF:311844.311845}. Their approach is based in cranio-facial development theory using geometrical ratios between different face regions to classify images into one of three age groups (babies, young adults and senior adults). They used frontal images in a very strict setup to be able to locate all face components. N. Ramanathan et al. \cite{1640784, Ramanathan2009131} used a similar approach in this case using 8 ratios rather than the 6 used by Y. Kwon et al.

The problem of this model is that it can only be applied to face images of young people in a growing age, since afterwards the facial geometry does not change as much. It is also a problem that the methods require frontal images, since it limits future applications.

\subsection{Active Appearance Models}\label{subsec:AAM}

\gls{aam} is a statistical shape model proposed by T.F Cootes et al. \cite{Cootes:2001:AAM:378040.378090}. This model contains the shape and grey-level appearance of the object of interest which can generalize to almost any valid example. This technique has been used to find the shape of faces by many researchers. A. Lanitis et al. \cite{791208} \cite{993553} \cite{Lanitis:2004:CDC:2225304.2226166} were the first to use the \gls{aam} model for age estimation by defining an ageing function $age=f(b)$, where $b$ is a vector containing the parameters learned by the \gls{aam}.

A. Lanitis et al. \cite{Lanitis:2004:CDC:2225304.2226166} also tried different classifiers such as Quadratic Functions, Shortest Distance Classifier, Supervised Neural Network and Unsupervised Neural Network. Among all of them they reported that Quadratic Functions where the ones performing best.

Later on, K. Luu et al. \cite{Luu:2009:AEU:1736406.1736456, LuuSSBS11}, use in \cite{Luu:2009:AEU:1736406.1736456} \gls{aam} with 68 facial landmarks to classify faces into young or adult classes and then use two specialised functions to finally determine the age. Luu's approach performs better than Lanitis' one \cite{Lanitis:2004:CDC:2225304.2226166}. In \cite{LuuSSBS11} they further improved the previous method by proposing Contourlet Appearance Models (CAM) which is more accurate and faster at calculating facial landmarks than \gls{aam}. This model has the ability of not only capturing global texture information like \gls{aam} but also local texture information using Nonsubsampled Contourlet Transform (NSCT) \cite{1703596}.

Chang et al. in \cite{5995437} used \gls{aam} model with particular ranking formulation of support vectors, OHRank. The approach uses cost-sensitive aggregation to estimate ordinal hyperplanes (OH) and ranks them according to the relative order of ages.

This model captures shape and texture information and in general performs better than the \textit{Anthropomorphic Models}. This method can deal with any range of ages rather than just with young ages like the previous model. However, as suggested by X. Geng et al. \cite{Geng:2006:LFA:1180639.1180711}, the ageing function is empirically determined, so there is no evidence suggesting that the relation between face and age is described just by a quadratic function.

\subsection{Ageing Pattern Subspace}\label{subsec:APS}

X. Geng et al. \cite{4359348, Geng:2006:LFA:1180639.1180711} were the first ones to explore \gls{ages} model. They define an \textit{ageing pattern} as a sequence of personal face images sorted in time order. Given a grey-scale face image $I$, where $I(x,y)$ determines the intensity of the pixel $(x,y)$, then an ageing pattern can be represented as a three-dimensional matrix $P$, where $P(x,y,t)$ is the intensity of the pixel $(x,y)$ in the face image at the time $t$. The images vector is filled with the available face images leaving empty the missing faces in the $t$ axis. Now, the images in the age pattern vector can be precessed and transformed into meaningful feature vectors.

In order to extract the features X. Geng use \gls{aam} as used in \cite{791208} since they capture the shape and texture of the face images. By representing ageing patterns in this way, the concepts of identity and time are naturally integrated into the data without any pre-assumptions.

The principal drawback of the \gls{ages} method is that assumes that there are images of the same individual at different ages, which is not true in many age databases. The databases that fulfil these requirements are small such as FG-NET \cite{FGNET} with just 1000 images and very few representations of the same individual over time.

\subsection{Age Manifold}\label{subsec:manifold}
 
Manifold learning methods are applied to find a sufficient embedding space and model the low-dimensional manifold data with a multiple linear regression function. Y. Fu et al. \cite{4523958, 4284917} were the first to propose a manifold embedding approach for the age estimation problem. 

The objective of this method is to find the low-dimensional representation in the embedded subspace capturing the intrinsic data distribution and geometric structure as well as its representation. G. Guo et al. \cite{Guo:2008:IHA:2319085.2321608} \cite{4531189} shows that the Orthogonal Locality Preserving Projections (OLPP) \cite{CHHZ06}  is a good an effective algorithm  to connect the manifold learning with subspace learning. In a posterior work \cite{5995404}, G. Guo et al. introduces a new approach, using kernel partial least square (KPLS) regression which reduces feature dimensionality and learn the ageing function in a single step.

Then T. Wu et al. \cite{journals/tifs/WuTC12} proposed to model the facial shapes as points on a Grassmann manifold. Age estimation is then considered as regression and classification problems on this manifold. Then, they proposed a method for combining this shape-based approach with other texture-based algorithms.

The main drawback of the age manifold representation is the large number of training instances required to learn the embedded manifold with statistical sufficiency.

\subsection{Appearance Models}\label{subsec:BIF}

Appearance Models focus on wrinkles, face texture and pattern analysis. From the beginning, researchers have tried to capture wrinkles and distinguish them from facial lines. Y. Kwon et al. \cite{Kwon:1999:ACF:311844.311845} proposed a wrinkle detector based on snakelets \cite{Kass88snakes:active} placed into key wrinkle areas of the face. Hayashi et al. \cite{969698} \cite{1195171} \cite{conf/icpr/HayashiYIK02} combined both shape and texture to estimate age and gender. In Hayashi's proposed approach, the skin is extracted based on a shape model and then a histogram equalization is applied to emphasize wrinkles.

Other researchers have used the texture descriptor \gls{lbp} \cite{Ahonen:2006:FDL:1175897.1176245} in the age estimation problem, such as \cite{4717926} \cite{6460367} obtaining good classification results with Nearest Neighbour and \gls{svm} classification algorithms. The Gabor \cite{Liu:2002:GFB:2319007.2320264} filter texture descriptor has also been used in the age estimation task  \cite{Gao:2009:FAC:1567988.1568003}, proving to be more discriminative than \gls{lbp}.

G. Guo et al. \cite{conf/cvpr/GuoMFH09} procposed to use \gls{bif} \cite{Riesenhuber99hierarchicalmodels} for age estimation via faces. The \gls{bif} descriptor tries to mimic how the visual cortex works, with a hierarchy of increasingly sophisticated representations. The \gls{bif} original model proposed by Riesenhuber and Poggio \cite{Riesenhuber99hierarchicalmodels} is based in a feed-forward model of the primate visual object recognition pathway, the ``HMAX'' model. The framework of the model contains alternative layers called simple (S) and complex (C) creating in each cycle a more elaborated representation. The S layers are created with a Gabor filtering on the input and the C layers generally operates a ``MAX'' operator over the previous S layer. G. Guo et al. \cite{conf/cvpr/GuoMFH09} modifies tha \gls{bif} model by changing the operator in the complex layer (C1) from ``MAX'' to ``STD''.

Han et al. \cite{han:age} used the \gls{bif} features in an hybrid classification framework improving the previous results with this descriptor. G. Guo et al. \cite{Guo2014761}, in a recent paper (2014), also used the \gls{bif} features, and focus to investigate a proposed single-step framework for joint estimation of age, gender and ethnicity. Both the Canonical Correlation Analysis (CCA) \cite{hotelling1936relations} and Partial Least Square (PLS) based methods were explored under the joint estimation framework.

In \cite{6553772}, Weng et al. employs a similar ranking technique than Chang et al. in \cite{5995437} called MFOR. \gls{lbp} histogram features are combined with principal components of \gls{bif}, shape and textural features of \gls{aam}, and \gls{pca} projection of the original image pixels. Fusion of texture and local appearance descriptors (\gls{lbp} and HOG features) have independently also been used for age estimation by Huerta et al. in \cite{HuFerPra14}.

The outstanding results obtained by some of these works point out the suitability of \gls{bif} features for the age estimation via faces task.

\subsection{Other Models}\label{subsec:oth}

Depending on the available data the age estimation problem changes. Many researchers have tackled the age estimation problem with different types of data. A short overview is described below.

N. Ramanathan et al. \cite{1709980} studied age progression of individual faces and proposed a method to perform face verification using a Bayesian age-difference classifier to improve the face verification algorithm.

A. Lanitis et al. \cite{5463396} collected a database of head and mouse tracking movements and then attempted to perform age estimation with this data. This preliminary study shows the potential of using this type of data.

Y. Mikihara et al. \cite{6117531} used a gait-based database to perform a viability study of age estimation using this type of data. The results show that in future research the combination of gait-based data and face-based age estimation could give very good results.

B. Xia et al. \cite{xia:hal-00904007} were the first to attempt age estimation using 3D face images. The obtained results show that the depth dimension has a very discriminative power in this problem. 

\section{Age Estimation Learning Algorithm}
Given an age representation, the next step is to determine the individual's age out of the ageing features. Age labels can be seen as a discrete set of classes or as a continuous label space, hence classification and regression learning methods can be used.

\subsection{Classification Methods}

The age estimation problem can be treated as a classification problem, where the solution space is discrete and the objective is to classify each face image into one of the age classes (a class could be an age range of several years or a single year).

A. Lanitis et al. \cite{Lanitis:2004:CDC:2225304.2226166} evaluated the performance of different classifiers such as quadratic function classifier, \gls{ann} and \gls{knn} classifier with their \gls{aam} model. Among the classification methods they tested, they claimed to perform better with the \gls{ann}, specifically the \gls{mlp}, obtaining 4.78 \gls{mae}. The authors also proposed some extensions, for example, training age specific classifiers in a hierarchical fashion. With the extended methods the authors reduced the error to 4.38 \gls{mae} with the \gls{mlp} and 3.82 \gls{mae} with the quadratic function (regression).

There have been previous proposals training neural networks,
which are able to learn complex mappings and deal with
outliers, for age estimation. In \cite{Lanitis:2004:CDC:2225304.2226166}, Lanitis et al. used \gls{aam} encoded face parameters as an input for the supervised training
of a neural network with a hidden layer. More recently, Geng
et al. \cite{6475129} tackled age estimation as a discrete classification problem using 70 classes, one for each age. The best algorithm proposed in this work (CPNN - Conditional Probability Neural Network) consists of a three-layered neural network, in which the input to the network includes both \gls{bif} features $x$ and a numerical value for age $y$, and the output neuron is a single value of the conditional probability density function $p(y|x)$. An extensive comparison of these classification schemes for age estimation has been reported in Fernandez et al. \cite{2014icprw}. In \cite{5995481}, Yang et al. used \gls{cnn} for age estimation under surveillance scenarios and recently Dong Yi et al. \cite{yiage} and Chenjing Yan et al. \cite{yan2014} have used \gls{cnn} in the age estimation problem reporting very promising results in the MORPH dataset (3.63 \gls{mae}). 

K. Ueki \cite{Ueki:2006:SAC:1126250.1126269} classified the images from the WIT\_DB database into 11 age groups using Gaussian models in a low-dimensional 2DLDA+LDA feature space using the EM Algorithm. The accuracy rates they achieved were $46.3\%$, $67.8\%$ and $78.1\%$ for age groups that were in the 5-year, 10-year and 15-year range respectively.

SVM have been also used for age classification, Guo et al. \cite{4544009} \cite{Guo:2008:IHA:2319085.2321608} trained an \gls{svm} for each pair of age classes and then using a binary tree search for testing, obtaining 5.55 \gls{mae} for females and 5.52 \gls{mae} for males in the YGA database and 7.16 \gls{mae} in the FG-NET database. 


\subsection{Regression Methods}

The age of an individual is nothing else than the time passed from the individual's birth, and time is a continuous dimension. Hence, the age estimation problem can be formulated as a regression problem where the objective is to find a regression function that explains the ageing in terms of the feature space.

A. Lanitis et al. \cite{Lanitis:2004:CDC:2225304.2226166} evaluated three regression functions, linear, quadratic and cubic and claimed that the quadratic function was the one which better described the age from their feature space. Y. Fu et al. \cite{4523958} \cite{4284917} used linear, quadratic and cubic regression functions as learning algorithms for the manifold age representation. As Lenitis showed, Fu also reported superior performance on the quadratic regression function, pointing out that cubic functions lead to over-fitting while linear functions lead to under-fitting.

Guo et al. \cite{4544009} \cite{Guo:2008:IHA:2319085.2321608} compared the performance of \gls{svr} against the Local Adjusted Robust Regression (LARR) performance in the YGA and the FG-NET databases, concluding that LARR performs a more accurate estimation, achieving 5.25 \gls{mae} in YGA for female images, 5.30 in YGA for male images and  5.07 in FG-NET database. In a posterior work \cite{conf/cvpr/GuoMFH09}, the authors improve the \gls{svr} performance in the FG-NET to 4.77 \gls{mae} by using \gls{bif} features.

\subsection{Hybrid Methods}

Many authors have proposed mixture frameworks, using both classification and regression learning algorithms. 

G. Guo et al. \cite{4563041} proposed a probabilistic fusion approach. They use Bayes' rule to derive the predictor and then a sequential fusion strategy, so the output of the regressor is used as an intermediate decision which is then fed to the classifier to aid or affect the decision space of the classifier. Their fusion approach has better performance than other single step methods which they compare with.

\gls{svm} and \gls{svr} where used by Han et al. \cite{han:age} in a hierarchical fashion. They proposed to use a binary decision tree with \glspl{svm} at each node to classify the images into different age ranges, which are coarsely assigned. Later, the age is fine grained by \glspl{svr} at the leaves. The \glspl{svr} are trained with 5-years-overlap between age ranges in order to reduce the misclassification error.


\begin{table}[h!]
	\centering
	\resizebox{\textwidth}{!}{
		\begin{tabular}{|l||c|M{3cm}|M{3cm}|M{3cm}|M{3.5cm}|N}
			%cline{1-7}
			\hline
			\textbf{Publication} & 
			\textbf{Year} &
			\textbf{Database (\#subjects, \#images)} &
			\textbf{Age Image Representation} &
			\textbf{Learning Method} & 
			\textbf{MAE} &\\[8pt] \hhline{=#=====}
			
			\multicolumn{1}{|p{2.5cm}||}{\textbf{A. Lanitis et al. \cite{993553}}} & 2002 & Private (60, 500) & AAM & Quadratic Aging Function & $3.94\pm3.8$&\\[5pt] \hline
			
			\multicolumn{1}{|p{2.5cm}||}{\textbf{A. Lanitis et al. \cite{Lanitis:2004:CDC:2225304.2226166}}} & 2004 & Private (40, 400) & AAM & Quadratic Aging Function & $3.82\pm5.58$&\\[5pt] \hline
			
			\multicolumn{1}{|p{2.5cm}||}{\textbf{X. Geng et al. \cite{Geng:2006:LFA:1180639.1180711}}} & 2006 & FG-NET & AGES & Regression & $6.77$ &\\[5pt] \hline
			
			\multicolumn{1}{|p{2.5cm}||}{\textbf{Y. Fu et al. \cite{4284917}}} & 2007 & YGA & Manifold & Linear regression function& $5\sim 6$ &\\[5pt] \hline
			
			\multicolumn{1}{|p{2.5cm}||}{\textbf{G. Guo et al. \cite{conf/cvpr/GuoMFH09}}} & 2009 & FG-NET, YGA & BIF & SVR & FG-NET/YGA $4.77$/F:$3.91$, M:$3.47$ &\\[5pt] \hline
			
			\multicolumn{1}{|p{2.5cm}||}{\textbf{K. Luu et al. \cite{Luu:2009:AEU:1736406.1736456}}} & 2009 & FG-NET & AAM & Hierarchical framework (SVM, SVR) & $4.37$ &\\[5pt] \hline
			
			\multicolumn{1}{|p{2.5cm}||}{\textbf{K. Luu et al. \cite{LuuSSBS11}}} & 2011 & FG-NET, PAL & Contourlet Appearance Model (CAM) & Hierarchical framework (SVM, SVR) & FG-NET/PAL $4.12$ / $6.0$ &\\[5pt] \hline
			
			\multicolumn{1}{|p{2.5cm}||}{\textbf{G. Guo et al. \cite{5995404}}} & 2011 & MORPH-II & BIF & Kernel PLS Regression & $4.18$ &\\[5pt] \hline
			
			\multicolumn{1}{|p{2.5cm}||}{\textbf{H. Han et al. \cite{han:age}}} & 2013 & FG-NET, MORPH-II, PCSO & BIF + Active Shape Models & Hierarchical framework (SVM, SVR) &  FG-NET/PCSO $4.7\pm24.8$ / $7.2\pm32.0$ &\\[5pt] \hline
			
			\multicolumn{1}{|p{2.5cm}||}{\textbf{G. Guo et al. \cite{Guo2014761}}} & 2014 & MORPH-II & BIF & rKCCA\footnote{Regularized Kernel Canonical Correlation Analysis} + Linear SVM & $3.92$ &\\[5pt] \hline
			
			\multicolumn{1}{|p{2.5cm}||}{\textbf{I. Huerta et al. \cite{HuFerPra14}}} & 2014 & FG-NET, MORPH-II & BIF + HOG + LBP + SURF + Gradient & rCCA\footnote{Regularized Canonical Correlation Analysis} & FG-NET/MORPH $4.17$ / $4.25$ &\\[5pt] \hline

		\end{tabular}
	}
	\caption[Summary of age estimation methods.]{Summary of some of the most important age estimation methods.}
	\label{tab:age-methods}
\end{table}

\section{Applications}
There are many real-world applications related to age estimation. Automatic age estimation is useful in situations where there is no need to specifically identify the individual, such as a government employee, but want to know his or her age.

\subsection{Security Control and Surveillance Monitoring}
In the last years security control and surveillance monitoring have gotten more relevant with the growth of internet content and the spread of technology that allows access to that content to under-age teenagers. Automatic age estimation systems can be used to prevent minors to buy alcohol in a grocery store, enter a bar or purchase tobacco from vending machines.

\subsection{Biometrics}
There are two types of biometric systems based on the number of traits used for recognition, unimodal biometric systems which consist on a single recognition trait and multimodal biometric systems, which combine evidences obtained from multiple sources \cite{MSU-CSE-99-39} such as fingerprints, iris, face, etc. The multimodal system is more robust, more reliable and secure against spoofing attacks. However, the data acquisition is much more troublesome than the unimodal. In order to overcome these inconveniences, soft-biometrics \cite{conf/icba/JainDN04}, such as age, hight, weigh, gender, ethnicity and eye colour, are used in combination with classic biometric traits. 

\subsection{Age-based Indexing Face Databases}
With the rise of interest for big data new and more efficient ways to retrieve data have to be developed. In large face image datasets, age can be used for indexing so that the queries to the dataset are simpler and faster. This is specially important in law enforcement where large image databases of suspects have to be filtered in order to find the most accurate suspects.

\subsection{Human-Computer Interaction and e-Commerce}
With the growth of e-commerce, companies want to offer a more personalized experience to their customers. Personalizing the offer or the product itself increases the user's satisfaction and the companies sells. Some examples of such policies are the following: Google \cite{Brin:1998:ALH:297810.297827} indexes the search results so the links that appear first appeal more to the user, Amazon \cite{Linden:2003:ARI:642462.642471} uses a recommender system to suggest products to the potential buyers according to their previous purchases, Netflix \cite{Koren:2009:MFT:1608565.1608614} held a competition in 2009 to create a film recommender system and gave a price of US \$1,000,000. Age estimation systems could have an important role in the sector since age is a discriminative feature for different client profiles. Visada \cite{visada} is an example of the use of age estimation for product recommendation.


\section{Age-based Datasets} \label{sec:ageDB}
There are many databases of faces in the literature, however, not so many capture the age of the individuals. This fact is due to the complexity of crawling such an information (if existent) from the usual fonts such as \textit{Flickr} or \textit{Facebook} and due to privacy issues. Moreover, the difficulty is even higher if the database contains chronological image series of individuals. Table \ref{tab:age-db} shows the most relevant databases used in the literature with number of samples, number of subjects, age range, type of age annotations and additional information if any. \textit{FG-NET} \cite{993553} is one of the first and most consolidated age databases, it is used to compare with other age estimation methods.

After an initial interest in automatic age estimation from images dated back to the early 2000s \cite{Lanitis:2004:CDC:2225304.2226166} \cite{993553} \cite{palDB}, research in the field has experienced a renewed interest from 2006 on due to the availability of large databases like \textit{MORPH-Album 2} \cite{1613043}, which contain 55 times more age-annotated images than the \textit{FG-NET} database.

\begin{table}[t!]

\centering
\resizebox{\textwidth}{!}{

\begin{tabular}{|l||c|c|c|c|M{2.5cm}|M{2.5cm}|M{2.5cm}|N}
	\hline
	\cline{1-8}
	\hline
	\textbf{Database} & \textbf{\#Faces} & \textbf{\#Subj.} & \textbf{Range} & \textbf{Type of age} & \textbf{Controlled Env.} & \textbf{Balanced age Distr.} & \textbf{Other annotation} &\\[8pt] \hhline{=#=======}
	
	\multicolumn{1}{|p{2cm}||}{\textbf{FG-NET \cite{993553, FGNET}}} & 1,002 & 82 & 0 - 69 & Real Age & No & No & 68 Facial Landmarks &\\[5pt] \hline
	
	\multicolumn{1}{|p{2cm}||}{\textbf{GROUPS \cite{gallagher_cvpr_09_groups}}} & 28,231 & 28,231 & 0 - 66+ & Age group & No & No & - &\\[5pt] \hline
	
	\multicolumn{1}{|p{2cm}||}{\textbf{PAL \cite{palDB}}} & 580 & 580 & 19 - 93 & Age group & No & No & - &\\[5pt] \hline
	
	\multicolumn{1}{|p{2cm}||}{\textbf{FRGC \cite{frgcDB}}} & 44,278 & 568 & 18 - 70 & Real Age & Partially & No & - &\\[5pt] \hline
	
	\multicolumn{1}{|p{2cm}||}{\textbf{MORPH2 \cite{1613043}}} & 55,134 & 13,618 & 16 - 77 & Real Age & Yes & No & - &\\[5pt] \hline
	
	\multicolumn{1}{|p{2cm}||}{\textbf{YGA \cite{4523958}}} & 8,000 & 1,600 & 0 - 93 & Real Age & No & No & - &\\[5pt] \hline
	
	\multicolumn{1}{|p{2cm}||}{\textbf{FERET\cite{Phillips1998295}}} & 14,126 & 1,199 & - & Real Age & Partially & No & - &\\[5pt] \hline
	
	\multicolumn{1}{|p{2cm}||}{\textbf{Iranian face \cite{4469272}}} & ~3,600 & 616 & 2 - 85 & Real Age & No & No & Kind of skin and cosmetic points \footnote{Surgical points, fracture or laceration on face.}  &\\[5pt] \hline
	
	\multicolumn{1}{|p{2cm}||}{\textbf{PIE \cite{1004130}}} & 41,638 & 68 & - & Real Age & Yes & No & - &\\[5pt] \hline
	
	\multicolumn{1}{|p{2cm}||}{\textbf{WIT-BD \cite{Ueki:2006:SAC:1126250.1126269}}} & 26,222 & 5,500 & 3 - 85 & Age group & No & No & - &\\[5pt] \hline
	
	\multicolumn{1}{|p{2cm}||}{\textbf{Caucasian Face Database \cite{burt1995perception}}} & 147 & - & 20 - 62 & Real Age & Yes & No & Shape represented in 208 key points &\\[5pt] \hline
	
	\multicolumn{1}{|p{2cm}||}{\textbf{LHI \cite{LHI}}} & 8,000 & 8,000 & 9 - 89 & Real Age & Yes & Yes & - &\\[5pt] \hline
	
	\multicolumn{1}{|p{2cm}||}{\textbf{HOIP \cite{HOIP}}} & 306,600 & 300 & 15 - 64 & Age Group & Yes & No & - &\\[5pt] \hline
	
	\multicolumn{1}{|p{2cm}||}{\textbf{Ni's Web-Collected Database \cite{Ni:2009:WIM:1631272.1631287}}} & 219,892 & - & 1 - 80 & Real Age & No & No & - &\\[5pt] \hline
	
	\multicolumn{1}{|p{2cm}||}{\textbf{OUI-Adience \cite{6906255}}} & 26.580 & 2.284 & 0 - 60+ & Age Group & No & No & Gender &\\[5pt] \hline

\end{tabular}

}

\caption{Age-based Databases and its characteristics.}
\label{tab:age-db}

\end{table}

\chapter{Data Collection} \label{chap:data}

As described in Section \ref{sec:ageDB} there are many age-based databases of facial images. However, all existing datasets are based on real age estimation. 

The idea of this work is to compare the performance between predicting real or apparent age labels. In order to do so, a web-based application has been developed using \textit{Facebook's API} to collect a database with these annotations.


\section{Web Application}

The aim of the web-based application was to speed up the collection and labelling processes and reach more people with broader backgrounds to create an age database as diverse as possible. These processes were implemented in a gamified \footnote{Gamification is the use of game thinking and game mechanics in non-game contexts to engage users in solving problems and increase users' self contributions. \cite{Deterding:2011:GDE:2181037.2181040}} fashion so the experience of the users with the application was satisfactory and engaging. 

The application uses the API of Facebook to create a ranking with the user's Facebook friends and add a factor of competitiveness to the game, and also to collect information about the labellers such as gender, age and nationality.

\begin{figure}[!h]
	\centering
	\begin{subfigure}[b]{0.3\textwidth}
		\includegraphics[width=\textwidth]{figures/age_app_1}
		\caption{Game Panel}
		\label{fig:game}
	\end{subfigure}%
	~ %add desired spacing between images, e. g. ~, \quad, \qquad, \hfill etc.
	%(or a blank line to force the subfigure onto a new line)
	\begin{subfigure}[b]{0.3\textwidth}
		\includegraphics[width=\textwidth]{figures/age_app_2}
		\caption{Gallery Panel}
		\label{fig:gallery}
	\end{subfigure}
	~ %add desired spacing between images, e. g. ~, \quad, \qquad, \hfill etc.
	%(or a blank line to force the subfigure onto a new line)
	\begin{subfigure}[b]{0.3\textwidth}
		\includegraphics[width=\textwidth]{figures/age_app_3}
		\caption{Ranking Panel}
		\label{fig:ranking}
	\end{subfigure}
	\caption{Age Recognition Application.} \label{fig:application}
\end{figure}

\subsection{Gamification Strategy}
The web-application is basically a platform to label and upload images, which is not funny or motivating. However, using gamification techniques the labeller engagement can be increased. The gamification strategies used were mainly three:

\begin{itemize}
	\item The users or players get \textbf{points} for uploading and voting (labelling) images. The closer the vote is to the apparent age (average current voted age) the more points the player gets. This strategy is pretended to persuade the users to wrongly label the images.
	\item Two \textbf{ranking} tables are shown to the users, one with the ranking positions of the users' Facebook friends and another showing the global classification. This strategy was created with the purpose of increase engagement between users, making them compete with each other. 
	\item A system of \textbf{achievements} was implemented, including the four types of achievements shown in the Table \ref{tab:achiev}.
	\begin{table}[!h]
		\centering
		\begin{tabular}{l|l}
			\textbf{Achievement} & \textbf{Description} \\ \hline
			Share it! &	Invite your friends to play and you will get stars.	\\
			Precision &	The better your guesses, the higher your rank. \\
			Vote! &	Vote on more images to get more stars.\\	
			Add Pictures! &	Upload more images to get more stars. \\
		\end{tabular}
		\caption{List of achievements with its description}
		\label{tab:achiev}
	\end{table}
	\item Thanks to one of the Challenge sponsors, California Naturel \footnote{\url{http://www.californianaturel.com/}}, a bunch of cosmetic lots were offer as a \textbf{prize} for the ones who lid the global ranking. The aim of the prize was to push participation further.
\end{itemize}


\subsection{Application Structure}
When the users access the web for first time they have to register using their Facebook account and accept the terms and conditions for the application usage (further explanation of the legal terms in the Subsection \ref{sec:trouble}). After registration the \textit{How to play} panel will be displayed giving a short description of all the parts of the application.

As it is shown in the Figure \ref{fig:application}, at the top of the site, there is a control menu that allows the users to move through all the site with just one click. A description of the different panels is written below:

\begin{itemize}
	\item \textbf{Profile}: In this section the users can keep track of their statistics (Number of uploaded images, number of votes, points obtained and global ranking position).
	\item \textbf{Game}: This is the main section (Figure \ref{fig:game}), four images are shown at the same time and the users are asked to guess the age of all of them. They can report the images if any of them is considered offensive, bad quality, more than one person appear in the image, etc.
	\item \textbf{Achievements}: In this panel the users can keep track of their achievements and shows the next goals.
	\item \textbf{Gallery}: In this section (Figure \ref{fig:gallery}) the players are able to see the images they have uploaded and see how many people have voted on them and get an estimation of the people opinion. They are also able to upload new images in this panel, while uploading they are asked to crop (if necessary) and specify the real age of the person in the image.
	\item \textbf{Ranking}:	The last panel (Figure \ref{fig:ranking}) allows users to compare their scores with their Facebook friends in the \textit{Friends Ranking} and to all the players in the \textit{Global Ranking}.
\end{itemize}

\subsection{Troubleshooting}\label{sec:trouble}
\note{Add}

We ask the users to upload images of a single person and we give them tools to crop the image if necessary, we also ask them to give the real age (or as close as possible) of the person in uploaded image, allowing more analysis and comparisons with real age and apparent age.


\section{HuPBA Age Dataset}

\note{Show statistic analysis of the database}

Few weeks after release the application we have already collected near 1000 images and near 10000 votes. These numbers will continue growing in order to generate the future competition. Some of the properties of the database which is being collected with the web application are listed below:

\begin{itemize}
	\item Thousands of faces labeled by many users.
	\item Images with background.
	\item Non-controlled environments.
	\item Non-labeled faces neither landmarks, making the estima-
	tion problem even harder.
	\item One of the first datasets in the literature including
	estimated age labeled by many users to define the ground truth
	with the objective of estimating the age.
	\item The evaluation metric will be pondered by the mean and
	the variance of the labeling by the participants.
	\item The dataset also provides for each image the real age
	although not used for recognition (just for analysis purposes).
\end{itemize}

In the same way for all the labelers we have their nationality, age, and gender, which will allow analyzing demographic and other interesting studies among the correlation of labelers. In relation to the properties of existing datasets shown in Table I, ours include labels of the real age of the individuals and the apparent age given by the collected votes, both age distributions are shown in the Figure 2. The images of our database has been taken under very different conditions, which makes it more challenging for recognition purposes.

\chapter{Method} \label{chap:method}
In this work tow methods combining state of the art techniques have been proposed. A full description is delineated in this Chapter.

\section{Preprocessing}

Given that the background information is not needed in any of the two proposed methods, it is required to first detect the faces from the images and second align all of them (Figure \ref{fig:preprocessing}).

\textbf{Face detection} was performed using the Viola Jones algorithm from the open source library OpenCV. A Haar Cascade classifier was trained with nearly 5000 face images and more than 9500 negative samples. After detecting the face, the image is cropped and resized into a $200 \times 200$.

\textbf{Face alignment}, 68 facial landmarks were extracted for every detected face with a version of the algorithm proposed by Shaoqing Ren et al. \cite{ren2014face}. The algorithm was trained with almost 5000 face labelled images from FG-NET, AFW, HELEN, IBUG and LFPW face databases. After extracting all the landmarks all the faces are aligned.

All the images where transformed into \textbf{grey scale} because of the variation of the data in terms of colour (some images are in grey scale other in RGB).

\begin{figure}[!h]
	\centering
	\includegraphics[width=\textwidth]{figures/preprocessing}
	\caption{Face Detection and Face Alignment}
	\label{fig:preprocessing}
\end{figure}


\section{Biologically Inspired Method}
As mentioned in Section \ref{subsec:BIF} Biologically Inspired Features (BIF) have been proven to work specially good in age estimation \cite{conf/cvpr/GuoMFH09}\cite{han:age}. In this work a method is proposed using BIF as age representation and a hierarchical framework with SVM and SVR as a learning method.
 
\subsection{Background}

Inspired by how the visual cortex works, Maximilian Riesenhuber and Tomaso Poggio \cite{Riesenhuber99hierarchicalmodels} proposed the BIF model for object recognition called ``HMAX'' based on the hierarchical model of the visual nervous system proposed by D. H. Hubel and T. N. Wiesel \cite{Hubel:62}. 

The ``HMAX'' model is composed by layers that will contain increasingly sophisticated representations. There are two type of layers called simple ($S_1$) and complex ($C_1$). Layer $S_1$ consists a battery of Gabor filters with different orientation and scales and layer $C_1$ is a pooling layer, in this model is a ``MAX'' pooling layer. There are many variations and extensions of this model made by many authors, Serre et al. \cite{4069258}\cite{1467551} introducing two more layers $S_2$ and $C_2$ or Mayers and Wolf \cite{Meyers:2008:UBI:1325290.1325298} uses a spacial $S_2$ layer called $S_2$ facial features (S2FF) for face recognition.

\begin{table}[!h]
	\centering
	\begin{tabular}{|M{2cm}|M{2cm}|M{1.5cm}|c|c|c|}
		\hline
		\multicolumn{3}{|c|}{$C_1$ Layer}                                               & \multicolumn{3}{c|}{$S_1$ Layer}                \\ \hline
		Scale Band $S$    & Pool grid $(N_S\times N_S)$    & Overlap $\Delta s$     & Filter size $s$  & Gabor $\sigma$ & Gabor $\lambda$ \\ \hline
		\multirow{2}{*}{Band 1} & \multirow{2}{*}{$6 \times 6$}   & \multirow{2}{*}{3}  & $5 \times 5$   & 2.0            & 2.5           \\
		&                                 &                     & $7 \times 7$   & 2.8            & 3.5           \\ \hline
		\multirow{2}{*}{Band 2} & \multirow{2}{*}{$8 \times 8$}   & \multirow{2}{*}{4}  & $9 \times 9$   & 3.6            & 4.6           \\
		&                                 &                     & $11 \times 11$ & 4.5            & 5.6           \\ \hline
		\multirow{2}{*}{Band 3} & \multirow{2}{*}{$10 \times 10$} & \multirow{2}{*}{5}  & $13 \times 13$ & 5.4            & 6.8           \\
		&                                 &                     & $15 \times 15$ & 6.3            & 7.9           \\ \hline
		\multirow{2}{*}{Band 4} & \multirow{2}{*}{$12 \times 12$} & \multirow{2}{*}{6}  & $17 \times 17$ & 7.3            & 9.1           \\
		&                                 &                     & $19 \times 19$ & 8.2            & 10.3          \\ \hline
		\multirow{2}{*}{Band 5} & \multirow{2}{*}{$14 \times 14$} & \multirow{2}{*}{7}  & $21 \times 21$ & 9.2            & 11.5          \\
		&                                 &                     & $23 \times 23$ & 10.2           & 12.7          \\ \hline
		\multirow{2}{*}{Band 6} & \multirow{2}{*}{$16 \times 16$} & \multirow{2}{*}{8}  & $25 \times 25$ & 11.3           & 14.1          \\
		&                                 &                     & $27 \times 27$ & 12.3           & 15.4          \\ \hline
		\multirow{2}{*}{Band 7} & \multirow{2}{*}{$18 \times 18$} & \multirow{2}{*}{9}  & $29 \times 29$ & 13.4           & 16.8          \\
		&                                 &                     & $31 \times 31$ & 14.6           & 18.2          \\ \hline
		\multirow{2}{*}{Band 8} & \multirow{2}{*}{$20 \times 20$} & \multirow{2}{*}{10} & $33 \times 33$ & 15.8           & 19.7          \\
		&                                 &                     & $35 \times 35$ & 17.0           & 21.2          \\ \hline
	\end{tabular}
	\caption{$S_1$ and $C_1$ parameters.}
	\label{tab:bif_param}
\end{table}

The proposed method in this work used the BIF model described by Guo et al. \cite{conf/cvpr/GuoMFH09} that uses ``STD'' pool operator instead of ``MAX'' operator (Figure \ref{fig:bif}). 

\begin{figure}[!h]
	\centering
	\includegraphics[width=\textwidth]{figures/BIF}
	\caption{Biologically Inspired Features (BIF)}
	\label{fig:bif}
\end{figure}


\textbf{$S_1$ layer} is formed by $S_1$ units which take as input a grey image $I$ of $200\times200$ pixels. These units are usually modelled by Gabor filters,

\begin{equation}
G(x,y) = exp(-\frac{X^2+\gamma^2 Y^2}{2\sigma^2})\cdot cos(\frac{2\pi}{\lambda}X)
\end{equation}

where $X=x\cos\theta + y\sin\theta$ and $Y = -x\sin\theta + y\cos\theta$ are the rotations of the filter with angle $\theta\in[0,\pi]$. The aspect ration $\gamma$ is set to $0.3$, the width $\sigma$, the wavelength $\lambda$ and the filter sizes $s$ are adjusted as in Table \ref{tab:bif_param}. These parameters are empirically determined based on reactions of the visual cortex to real stimuli \cite{4069258}. 

\textbf{$C_1$ layer} is formed by $C_1$ units which pool over the $S_1$ units with the same orientation and from the same scale band (see Table \ref{tab:bif_param}). Each scale band $S$ contains two adjacent filter sizes, for instance, scale band 1 contains filter with sizes $5\times5$ and $7\times7$. The scale band also determines the sizes of the neighbourhood over which the $C_1$ units pool ($N_S\times N_S$). The pooling operator used in this method is the ``STD'' operator proposed by Guo et al. \cite{conf/cvpr/GuoMFH09},

\begin{equation}
std_{j,j+1} = \sqrt{\frac{1}{N_S\times N_S}\sum_{i=1}^{N_S\times N_S}(F_i - \bar{F})^2} 
\end{equation}

where $F_i$ is the maximum value of two consecutive $S_1$ units output in the same scale band at pixel index $i$, 

\begin{equation}
F_i = \max(x_i^j,x_i^{j+1})
\end{equation}

where $x_i^j$ and $x_i^{j+1}$ are the filtered values with scales $j$ and $j+1$ at position $i$. $\bar{F}$ is the mean of the filtered values in the neighbourhood $N_S\times N_S$.

The $C_1$ units are concatenated into a single feature vector for each input image.


\subsection{System overview}

This method combines a hierarchical framework and a hybrid face model, mixing BIF features and face shape landmarks to optimize the performance. The pipeline consists of three steps (see Figure \ref{fig:pipeline}).

\begin{figure}[!h]
	\centering
	\includegraphics[width=\textwidth]{figures/pipeline}
	\caption{Biologically Inspired Method Pipeline}
	\label{fig:pipeline}
\end{figure}

\textbf{Feature extraction}: 
After the face alignment is done, BIF features are extracted from each face image. Because of the high dimensionality of the data ($\sim3000$ BIF features and 138 shape landmark coordinates) feature reduction is needed. Principal Component Analysis (PCA) is used to reduce dimensionality of both shape and BIF. As shows Figure \ref{fig:pipeline} BIF and shape features are concatenated into a single feature vector.

\textbf{Age Group Classification}:
In this stage the images are classified into three age groups with age ranges (0-18), (19-45) and (45-100). The classifier used is a linear SVM as suggested in \cite{4531189}.

\textbf{Age Estimation}:
Three SVR with RBF kernel were trained to perform age regression in a specific age range. The data used to fit the regressors were images within 5-years-overlap between age ranges to reduce the misclassification error as proposed in \cite{han:age}. 

In order to improve the system performance a last fusion step is done by averaging the two nearest age regressors, i.e. if the age group classifier predicts that a given face belongs to the adult age group (19-45)] and the adult-specific regression function predicts that the exact age is 43, then in the fusion step would average the adult regressor with the elder regressor.


\section{Deep Learning Method}
\subsection{Background}
\subsection{System overview}
\chapter{Results} \label{chap:experiments}

This Chapter describes the experiments done in this work and the results achieved.

\section{Datasets}
Two datasets were used in the experiments, our collected dataset HuPBA-AgeGuess and the classic benchmark database FG-NET.

\textbf{FG-NET} consists of 1002 frontal face images of 82 different individuals. The image quality varies a lot in the dataset since there are images in grey scale and RGB. The face position is frontal and under similar illumination conditions. The dataset also contain 68 facial landmarks for each face image.

\begin{figure}[!h]
	\centering
	\includegraphics[width=\textwidth]{figures/FGNET_sample}
	\caption{FG-NET image samples.}
	\label{fig:imgSample1}
\end{figure}

\textbf{HuPBA-AgeGuess} dataset (\ref{fig:imgSample2}) used in this project is a subset of the 4865 images filtered out by a minimum number of votes per image. This subset contains 3398 face images. The images are captured in the wild so the faces position vary up to $\pm90º$ and the illumination is also different in every picture.

\begin{figure}[!h]
	\centering
	\includegraphics[width=\textwidth]{figures/HuPBA_sample}
	\caption{HuPBA-AgeGuess image samples.}
	\label{fig:imgSample2}
\end{figure}

Given the characteristics of the used descriptors in this work, both databases could be incremented by computing the mirror image, doubling the number of faces of each dataset (i.e. 2004 faces for the FG-NET and 6796 face images for the HuPBA-AgeGuess dataset).

\section{Evaluation Metrics} 

The two most commonly used evaluation metrics in age estimation are \acrfull{mae} and \gls{cs}.

The \gls{mae} is described as,

\begin{equation}
MAE = \frac{1}{n_{s}}\sum_{i=0}^{n_{s}-1} |e_i|,
\end{equation}

where $n_s$ is the number of samples and $e_i$ is the error of the $ith$ instance, i.e. $e_i = |\hat{y_i} - y_i|$ where $y_i $ is the real label and $\hat{y_i}$ is the predicted label. This metric tells the average number of years that the prediction is wrong.

\gls{cs} is defined as the percentage of test images such that the absolute error is not higher than a threshold, $t$ (in years). i.e.,

\begin{equation}
CS(t) = (1 - \frac{1}{n_{s}}\sum_{i=0}^{n_{s}-1} h(|\hat{y_i} - y_i| - t))\cdot 100
\end{equation}
\begin{equation}
h(x) = 
\begin{cases}
1,				& \text{if } x \geq 0\\
0,              & \text{otherwise},
\end{cases}
\end{equation}

where $y_i$ is the age label of the $ith$ test image and $\hat{y_i}$ is the age prediction of the $ith$ test image.
 

\section{Experimental Settings}
This section describes the experimental setup and the parameters used in the two proposed methods.

\subsection{Biologically Inspired Method}

As described in Section \ref{sec:BIF} a \gls{svm} classifier and three \gls{svr} regressors were trained for this method. 

In order to find the best parameters a grid search was performed to find the best parameters. The parameters that formed the search space were the ones required by the \gls{svm} and \glspl{svr} and its kernels, which are the following:

\begin{itemize}
	\item \textbf{Penalty term ($C$)}: This is the Support Vector parameter that deals with the cost of a misclassification over all the classification task. 
	
	The $C$ parameters tried in the search space were between 0.1 and 2. The best parameters in the Age Group Classification were between 1 and 2 and the best parameters for the \glspl{svr} were between 0.1 and 1.
	
	\item \textbf{Influence term ($\gamma$)}: It is the \gls{rbf} kernel parameter. It determines how far the influence of a single training example reaches.
	
	The $\gamma$ parameters used were between 0.001 and 1, being between 0.001 and 0.01 the ones with better performance.
\end{itemize}

In order to train and validate the parameters a nested 10-fold cross validation technique was used.

To evaluate the performance of the age group classifier it is needed to take into account that the labels are a disjoint segmentation of the spectrum of possible ages. Therefore, the in the performance evaluation the problem is treated as a regression problem rather than a classification, qualifying the predictions based on how far they are from the real label. The error measure chose is the \gls{mae}, normalizing the labels by the number of classes first.

\subsection{Deep Learning Method}

The FG-NET dataset has only 1002, clearly not enough to train the described deep \gls{cnn} in Section \ref{sec:deep}. Therefore, this method has only been tested with the HuPBA-AgeGuess database.

Depending on the dataset the experimental setting were different. In order to validate and test the method a 10-fold cross validation was performed splitting the data into $80\%$ training set, $10\%$ validation set and $10\%$ testing set. The data was split into batches each of them containing 213 instances. In this way each fold will contain an exact number of batches, 3. The error function used during the training was the mean square error. The network also needs the following parameters to be set,

\begin{itemize}
	\item \textbf{Learning Rate ($\eta $)}: The learning rate is used in the backpropagation stage to modify the net weights.
	
	The learning rate used in the experiments was set to 0.002. 
	
	\item \textbf{Epoch}: In training a neural network, the term epoch is used to describe a complete pass through all of the training patterns. The error is reduced over the epoch since the weights are updated within each epoch.
	
	The number of epochs used in the training of the proposed \gls{cnn} were 400.
\end{itemize}

\section{Analysis of the Experiments}

% además de la discusión de los resultados recuerda poner muchas imágenes:

%-ejemplos de landmark fitting (1 página de imágenes)

%-imagenes ejem de perfecto reconocim- algunas que fallen.

%-en el caso de deep learning estaría bien mostrar las features encontradas!!!!! (los filtros) e incluso el resultado de aplicar los filtros a imagenes de caras !!!!

%incluir discusion de pros contras y comparativa de los 2 métodos.

%cuando lo tratas como multi-clasificación mostrar visualmente con colores la matriz de confusión

In this section the results of all the experiments is described and analysed in detail. The experiments were run in the two datasets, FG-NET and HuPBA-AgeGuess, amd with the two different age labels (just in HuPBA-AgeGuess).

\subsection{Face Detection and Alignment}

This preprocessing step was just applied to the HuPBA-AgeGuess dataset since the FG-NET already contains the 68 facial landmarks manually placed. The database HuPBA-AgeGuess was preprocessed as described in Section \ref{sec:preproces}. The face detection was $92.52\%$ accurate and face alignment was $84.54\%$ accurate, hence the total accuracy of the preprocessing was $78.26\%$. Some examples of good and bad face detections and alignments are shown in Figure \ref{fig:sampPP}. The detection and alignment were tricky since the images are not taken in a controlled environment, the multiple illumination changes and the multiple head orientations made the preprocessing very difficult.

\begin{figure}[!h]
	\centering
	\begin{subfigure}[b]{\textwidth}
		\includegraphics[width=\textwidth]{figures/good_pp}
		\caption{Well detected and aligned faces.}
		\label{fig:goodPP}
	\end{subfigure}%
	
	 %add desired spacing between images, e. g. ~, \quad, \qquad, \hfill etc.
	%(or a blank line to force the subfigure onto a new line)
	\begin{subfigure}[b]{\textwidth}
		\includegraphics[width=\textwidth]{figures/bad_detection}
		\caption{Samples of bad face detection.}
		\label{fig:badDet}
	\end{subfigure}
	
	 %add desired spacing between images, e. g. ~, \quad, \qquad, \hfill etc.
	%(or a blank line to force the subfigure onto a new line)
	\begin{subfigure}[b]{\textwidth}
		\includegraphics[width=\textwidth]{figures/bad_alignment}
		\caption{Samples of bad face alignment.}
		\label{fig:badAl}
	\end{subfigure}
	\caption{Samples of preprocessed faces.}\label{fig:sampPP}
\end{figure}


\subsection{Biologically Inspired Method}

The \gls{mae} achieved by this method in both databases is shown in the Table \ref{tab:BIFresults}. As it was expected the algorithm performed better with the FG-NET dataset since the landmarks were manually placed. Between both experiments with the HuPBA-AgeGuess dataset this method performed best in the apparent age labels.

\begin{table}[!h]
	\centering
	\begin{tabular}{|l||c|}
		\hline
		\multicolumn{1}{|c||}{\textbf{Database}} & \textbf{\gls{mae}}\\ \hhline{=#=}
		FG-NET & $7.99$\\ 		\hline
		HuPBA-AgeGuess real age & $10.73$\\ \hline
		HuPBA-AgeGuess apparent age & $9.34$\\ \hline
			
	\end{tabular}
	\caption{\acrshort{mae} achieved with \acrshort{bif}.}
	\label{tab:BIFresults}
\end{table}

Figure \ref{fig:cumS_BIF} shows the \gls{cs} obtained with this method. It shows that the performance is better over apparent age than real age.

\begin{figure}[!h]
	\centering
	\includegraphics[width=\textwidth]{figures/cum_score}
	\caption{\acrshort{cs} achieved with the Biologically Inspired Method.}
	\label{fig:cumS_BIF}
\end{figure}

The group classification achieves $0.895$ accuracy with FG-NET and $0.87$ with HuPBA-AgeGuess dataset. Figure \ref{fig:cm} shows the confusion matrix obtained in each one of the three datasets. In the FG-NET database there is confusion between the adult and the elder faces, and in the HuPBA-AgeGuess young and elder are confused with adults. This bad result is product of different factors, the high number of face instances with age between 18 and 45 years old bias the classifier towards the adult class.

\begin{figure}[!h]
	\centering
	\begin{subfigure}[b]{0.5\textwidth}
		\includegraphics[width=\textwidth]{figures/FG-NETconf_mat}
		\caption{FG-NET dataset.}
		\label{fig:cmfgnet}
	\end{subfigure}%
	~
	%add desired spacing between images, e. g. ~, \quad, \qquad, \hfill etc.
	%(or a blank line to force the subfigure onto a new line)
	\begin{subfigure}[b]{0.5\textwidth}
		\includegraphics[width=\textwidth]{figures/real_conf_mat}
		\captionsetup{width=1.1\textwidth}
		\caption{HuPBA-AgeGuess with Real Age labels.}
		\label{fig:cmreal}
	\end{subfigure} %
	
	\vspace{0.5cm}
	%add desired spacing between images, e. g. ~, \quad, \qquad, \hfill etc.n
	%(or a blank line to force the subfigure onto a new line)
	\begin{subfigure}[b]{0.5\textwidth}
		\includegraphics[width=\textwidth]{figures/apparent_conf_mat}
		\caption{HuPBA-AgeGuess with Apparent Age labels.}
		\label{fig:cmapp}
	\end{subfigure}
	\caption{Confusion Matrices of the Age Group Classification.}\label{fig:cm}
\end{figure}

The poor age group classification causes an overall bad performance of the age estimation algorithm with the HuPBA-AgeGuess dataset.

\note{Study the CS and MAE by group of ages (maybe images in Appendix)}

\subsection{Deep Learning Method}

The Table \ref{tab:CNNresults} shows the achieved \gls{mae} with this method with the HuPBA-AgeGuess database and both label annotations. The \gls{mae} of the method predicting real age is $40.77\%$ higher than predicting apparent age.

\begin{table}[!h]
	\centering
	\begin{tabular}{|l||c|}
		\hline
		\multicolumn{1}{|c||}{\textbf{Database}} & \textbf{\gls{mae}}\\ \hhline{=#=}
		FG-NET & - \\ 		\hline
		HuPBA-AgeGuess real age & $10.29$\\ \hline
		HuPBA-AgeGuess apparent age & $8.71$\\ \hline
		
	\end{tabular}
	\caption{\acrshort{mae} achieved with \acrshort{cnn}.}
	\label{tab:CNNresults}
\end{table}

Figure \ref{fig:cumS_CNN} shows the \gls{cs} obtained with this method. It shows that the performance is better over apparent age than real age.

\begin{figure}[!h]
	\centering
	\includegraphics[width=\textwidth]{figures/cum_score}
	\caption{\acrshort{cs} achieved with the \acrshort{cnn} Method.}
	\label{fig:cumS_CNN}
\end{figure}

\note{Show weight masks (maybe in the Apendix)}

\subsection{Comparison}

From the two tables of results (Tables \ref{tab:BIFresults} and \ref{tab:CNNresults}) can be seen that the \gls{cnn} method achieves better results with the HuPBA-AgeGuess, reducing the \gls{mae} $0.65$ years with apparent age and $0.44$ years with real age annotations.

It also can be observed that the performance with apparent age is in all the cases higher than with real age.

Despite the fact that the FG-NET database could not be tested with the \gls{cnn} method, the lowest \gls{mae} achieved in the experiments was using the 

In both \gls{cnn} and \gls{bif} methods the FG-NET dataset achieved better results. This is because the images in that dataset are frontal faces and the 68 landmarks were manually placed, while the HuPBA-AgeGuess database contains non frontal face images and the landmarks are automatically regressed.

%\begin{table}[!h]
%	\centering
%	\begin{tabular}{|l|l|c|}
%		\hline
%		\multicolumn{1}{|c|}{\textbf{Method}} & 
%		\multicolumn{1}{|c|}{\textbf{Database}} & \textbf{\gls{mae}}\\ \hline\hline%\hhline{====}
%		\multirow{3}{*}{Bio-Inspired Method} & FG-NET & $7.99$\\ 		\cline{2-3} 
%		& HuPBA-AgeGuess real age & $10.73$\\ \cline{2-3}
%		& HuPBA-AgeGuess apparent age & $9.34$\\ \hline\hline%\hhline{====}
%		\multirow{3}{*}{CNN Method} & FG-NET & \\ \cline{2-3}
%		& HuPBA-AgeGuess real age & $10.29$\\ \cline{2-3}
		
%		& HuPBA-AgeGuess apparent age & $8.71$ \\ \hline
		
		
%	\end{tabular}
%	\caption{\acrshort{mae} of the different experiments.}
%	\label{tab:results}
%\end{table}
\chapter{Conclusions and Future work} \label{chap:conclusions}

% References
\renewcommand{\bibname}{\refname} % To change the heading of the chapter from Bibliography to References
\addcontentsline{toc}{chapter}{\bibname}
\bibliographystyle{plain}
%\nocite{*} % Should be removed at the end to make sure all references are cited
\bibliography{references}

%\glsaddall
\printglossaries

\end{document}


%1. Introduction
%  1.1. Motivation
%  1.2. Goals
%2. State of the art
%3. Resources
%  3.1. NAO and Wifibot
%  3.2. Kinect v2
%  3.3. ROS & SMACH
%  3.4. SOAR?
%4. The System
%  4.1. Architecture/Overview -> Diagramas, tipos de gestos
%  4.2. Computer Vision: Real Time OnLine Gesture Recognition
%            4.2.1. Dynamic Time Warping
%            4.2.2. Skeletal extracted features per gesture
%  4.3. Mobile Robotics: (Cognitive) Interaction
%            4.3.1. State Machines
%            4.3.2. Robot interaction (en este apartado 4.3 . no tengo muy claro como segmentarlo y que poner aún...)
%            4.3.3. Moving to a more general, cognitive approach (explicar cómo he jugado con SOAR y cómo se podria aplicar)
%5. Experimental evaluation
%            5.1 Data, methods, and settings
%            5.2 Gesture recognition evaluation
%            5.3 User experience evaluation (preguntas, experimento, y opiniones, todo junto)
%6. Conclusions and future work
%References